\noindent\textbf{Elección del procedimiento.} Para excitación armónica, el módulo de la respuesta normalizada es $M(r)=[(1-r^2)^2+(2\zeta r)^2]^{-1/2}$. Se evalúa por separado en cada razón de frecuencias. El caso $r=1$ es coincidencia con la frecuencia natural, no exactamente el máximo con amortiguamiento finito.

\begin{enumerate}

\item En $r=1$ se anula la parte elástica del denominador.

\[\fitmath{M_{1}=\frac{1}{2\cdot \zeta}}\]

\[\fitmath{M_{1}=\frac{1}{2\cdot (0{,}05)}\approx 10\,}\]

\item Conservar tanto la parte real como la disipativa para $r=0{,}5$.

\[\fitmath{\mathrm{Mhalf}=\frac{1}{\sqrt{{\left(1-{0.5}^{2}\right)}^{2}+{\left(2\cdot \zeta\cdot 0{,}5\right)}^{2}}}}\]

\[\fitmath{\mathrm{Mhalf}=\frac{1}{\sqrt{{\left(1-{0.5}^{2}\right)}^{2}+{\left(2\cdot (0{,}05)\cdot 0{,}5\right)}^{2}}}\approx 1{,}33038\,}\]

\item Evaluar la condición pedida $r=0{,}9$.

\[\fitmath{\mathrm{Mr}=\frac{1}{\sqrt{{\left(1-{r}^{2}\right)}^{2}+{\left(2\cdot \zeta\cdot r\right)}^{2}}}}\]

\[\fitmath{\mathrm{Mr}=\frac{1}{\sqrt{{\left(1-{(0{,}9)}^{2}\right)}^{2}+{\left(2\cdot (0{,}05)\cdot (0{,}9)\right)}^{2}}}\approx 4{,}756515\,}\]

\end{enumerate}
